
% ═══════════════════════════════════════════════════════════════════════
% THEORETICAL ANALYSIS — CompTS Safety and Regret Guarantees  
% ═══════════════════════════════════════════════════════════════════════

\subsection{Problem Formulation}

We formulate compliance-aware query routing as a conservative contextual 
bandit. At each step $t = 1, \ldots, T$, a query $q_t$ arrives with 
unknown sensitivity $s_t \in \{0, 1, 2, 3\}$ corresponding to our 
four-tier taxonomy. A classifier produces estimate $\hat{s}_t$ with 
confidence $\gamma_t$. The router selects platform $a_t$ from action 
set $\mathcal{A} = \{a_1, \ldots, a_m\}$, each with clearance level 
$L(a_j)$ and cost $c_j$. A \emph{compliance violation} occurs when 
$L(a_t) < s_t$, i.e., a query is routed to a platform lacking 
sufficient clearance.

\begin{definition}[Compliance-Constrained Routing]
A routing policy $\pi$ is $\varepsilon$-compliant if for all $t$:
\begin{equation}
\Pr[L(\pi(q_t)) < s_t] \leq \varepsilon
\end{equation}
The goal is to find the $\varepsilon$-compliant policy minimizing 
cumulative cost $\sum_{t=1}^{T} c(\pi(q_t))$.
\end{definition}

\subsection{CompTS Algorithm}

At each step, CompTS maintains a posterior over $s_t$ using the 
classifier softmax output as likelihood and a Dirichlet prior updated 
from routing history. The safe action set is:
\begin{equation}
\mathcal{A}_{\text{safe}}(t) = \{a_j \in \mathcal{A} : 
\Pr_{\text{post}}[L(a_j) < s_t] \leq \varepsilon\}
\end{equation}
CompTS selects the cost-minimizing action from $\mathcal{A}_{\text{safe}}(t)$ 
with Thompson Sampling exploration.

% ── Theorem 1 ────────────────────────────────────────────────────────

\begin{theorem}[Per-Step Safety Guarantee]
\label{thm:safety}
Let $f$ be a classifier with tier-3 recall $r$ and calibration error 
$\alpha$, and let CompTS use confidence threshold $\tau$ and violation 
threshold $\varepsilon$. Then for each step $t$:
\begin{equation}
\Pr[\text{violation at } t] \leq p_{\phi} \cdot (1-r) \cdot p_\tau \cdot \varepsilon
\end{equation}
where $p_\phi = \Pr[s_t = 3]$ is the base rate of restricted queries 
and $p_\tau = \Pr[\gamma_t \geq \tau \mid s_t = 3, \hat{s}_t \neq 3] 
\leq \alpha$ for calibrated classifiers.
\end{theorem}

\begin{proof}
A violation at step $t$ requires the conjunction of four events:
\begin{align}
E_1 &: s_t = 3 \quad &&\text{(query contains PHI)} \\
E_2 &: \hat{s}_t < 3 \quad &&\text{(classifier misclassifies)} \\
E_3 &: \gamma_t \geq \tau \quad &&\text{(confidence exceeds fallback)} \\
E_4 &: s_t \notin \mathcal{S}_{\text{safe}}(a_t) \quad &&\text{(posterior check fails)}
\end{align}

\textbf{Bounding $E_1$:} $\Pr[E_1] = p_\phi$ by the data distribution.

\textbf{Bounding $E_2 | E_1$:} Given $s_t = 3$, the classifier 
misclassifies with probability $1 - r$, where $r$ is the T3 recall.

\textbf{Bounding $E_3 | E_1 \cap E_2$:} Given a T3 query is 
misclassified, the confidence $\gamma_t$ exceeds $\tau$ with 
probability $p_\tau$. For a calibrated classifier (ECE $\leq \alpha$), 
a prediction with confidence $\gamma_t \geq \tau$ is correct with 
probability at least $\tau - \alpha$. Thus, the probability of being 
both wrong and confident is bounded by $p_\tau \leq \alpha / (1-r)$ 
for calibrated models. Conservatively, $p_\tau \leq \alpha$.

\textbf{Bounding $E_4 | E_1 \cap E_2 \cap E_3$:} Given all three 
preceding events, CompTS selects $a_t \in \mathcal{A}_{\text{safe}}(t)$, 
where by construction:
$\Pr_{\text{post}}[L(a_t) < s_t] \leq \varepsilon$. Thus 
$\Pr[E_4 \mid E_1 \cap E_2 \cap E_3] \leq \varepsilon$.

By the chain rule:
\begin{equation}
\Pr[\text{viol}_t] = \Pr[E_1]\Pr[E_2|E_1]\Pr[E_3|E_1,E_2]\Pr[E_4|E_1,E_2,E_3]
\leq p_\phi (1-r) p_\tau \varepsilon
\end{equation}
\end{proof}

\begin{corollary}[Global Safety]
Over $T$ queries, the expected number of violations satisfies:
$\mathbb{E}[\text{violations}] \leq T \cdot p_\phi (1-r) p_\tau \varepsilon$.
By Markov's inequality, $\Pr[\text{violations} \geq 1] \leq T \cdot p_\phi(1-r)p_\tau\varepsilon$.
Setting $\varepsilon = \delta / (T \cdot p_\phi(1-r)p_\tau)$ yields 
$\Pr[\text{any violation}] \leq \delta$.
\end{corollary}

% ── Theorem 2 ────────────────────────────────────────────────────────

\begin{theorem}[Regret Bound]
\label{thm:regret}
Let $\pi^*$ be the optimal $\varepsilon$-compliant policy. The 
cumulative regret of CompTS is:
\begin{equation}
R(T) = \sum_{t=1}^{T}\left[c(\text{CompTS}_t) - c(\pi^*_t)\right]
\leq C_1 d\sqrt{T\log T} + C_2 T \cdot p_{\text{low}} \cdot \Delta_c
\end{equation}
where $d$ is the context embedding dimension, 
$p_{\text{low}} = \Pr[\gamma_t < \tau]$ is the fallback rate, and 
$\Delta_c = \max_j c_j - \min_j c_j$ is the cost range.
\end{theorem}

\begin{proof}[Proof sketch]
The regret decomposes into two components:

\textbf{Exploration regret:} At each step where $\gamma_t \geq \tau$, 
CompTS performs Thompson Sampling over the safe action set 
$\mathcal{A}_{\text{safe}}(t)$. Since $\pi^*$ also satisfies 
compliance constraints, $\pi^*(q_t) \in \mathcal{A}_{\text{safe}}(t)$ 
for all $t$. The constrained action space is a subset of $\mathcal{A}$, 
so by the standard Thompson Sampling regret bound 
\citep{agrawal2013thompson}: $R_{\text{explore}}(T) \leq C_1 d\sqrt{T\log T}$.

\textbf{Conservatism penalty:} When $\gamma_t < \tau$, CompTS falls 
back to the maximum-clearance platform, adding excess cost 
$\Delta_c$ per fallback query. The expected penalty is:
$R_{\text{conserv}}(T) \leq T \cdot p_{\text{low}} \cdot \Delta_c$.

Total regret: $R(T) \leq R_{\text{explore}}(T) + R_{\text{conserv}}(T)$.
\end{proof}

% ── Theorem 3 ────────────────────────────────────────────────────────

\begin{theorem}[Ensemble Amplification]
\label{thm:ensemble}
With $k$ independent classifiers each having T3 recall $r > 0.5$, 
using majority vote:
\begin{equation}
\Pr[\text{violation}] \leq \varepsilon \cdot \sum_{i=\lceil k/2\rceil}^{k}
\binom{k}{i}(1-r)^i r^{k-i}
\end{equation}
which decreases exponentially:
$\Pr[\text{violation}] \leq \varepsilon \cdot \exp(-k \cdot D(1/2 \| 1-r))$
where $D(\cdot\|\cdot)$ is the KL divergence.
\end{theorem}

\begin{proof}
With $k$ independent classifiers, the majority vote correctly 
identifies a T3 query when at least $\lceil k/2 \rceil$ classifiers 
are correct. The recall of the majority vote ensemble is:
\begin{equation}
r_{\text{ens}} = \sum_{i=\lceil k/2 \rceil}^{k} \binom{k}{i} r^i (1-r)^{k-i}
= 1 - \sum_{i=\lceil k/2\rceil}^{k}\binom{k}{i}(1-r)^i r^{k-i}
\end{equation}
Substituting $r_{\text{ens}}$ into Theorem~\ref{thm:safety} 
(replacing $r$ with $r_{\text{ens}}$):
\begin{equation}
\Pr[\text{viol}] \leq p_\phi(1-r_{\text{ens}})p_\tau\varepsilon
= p_\phi \left[\sum_{i=\lceil k/2\rceil}^{k}\binom{k}{i}(1-r)^i r^{k-i}\right] p_\tau \varepsilon
\end{equation}
The bracketed term is the tail of a Binomial$(k, 1-r)$ distribution 
above $k/2$. By the Chernoff bound, for $r > 0.5$:
$\sum_{i=\lceil k/2\rceil}^{k}\binom{k}{i}(1-r)^i r^{k-i} 
\leq \exp(-k \cdot D(1/2 \| 1-r))$
\end{proof}

\begin{remark}
For $r = 0.97$ and $k = 5$: $\Pr[\text{majority wrong}] \approx 
\binom{5}{3}(0.03)^3(0.97)^2 + \binom{5}{4}(0.03)^4(0.97)^1 + 
(0.03)^5 \approx 2.55 \times 10^{-4}$, a $100\times$ improvement 
over single-classifier violation probability.
\end{remark}
